\section{Cote 134-8: higher coproduct surjectivity is not forced by the sheet}
\label{sec:arch134}

\subsection{Sheet and inherited statements}
The tensor, scalar adjunctions, global normalization, and interval conditions
are at cote~134-8, physical p.~48, footer~47,
\texttt{ARCHIVE\_134-8\_scan048.tex}; the displayed \emph{Thm(?)} is on
physical p.~49, footer~48, \texttt{ARCHIVE\_134-8\_scan049.tex}.
We retain the interpretation of those predicates fixed in
ALPOGE\_BRIDGES, BR-13403--06 and F-SM-32--44.  In particular the
scalar embedding is exact and has both adjoints; the tensor is symmetric,
associative, unital and bilinear; and the interval map remains monic after
tensoring with every object and has fully derived-acyclic target.
The construction below has a tensor preserving all colimits in each
variable, and its interval and interval quotient are flat.  Thus it does
not exploit a weaker interpretation of bilinearity.

F-SM-44 establishes that the canonical comparison
\[
 b_n:\bigoplus_s L_n\varphi(X_s)\longrightarrow
                  L_n\varphi\!\left(\bigoplus_sX_s\right)
\]
is monic in every degree and an isomorphism in degrees zero and one.
The new result is that its higher surjectivity can fail under \emph{all}
the sheet's tensor and interval hypotheses.  The localization conclusion
nevertheless holds in this example.  Consequently the higher-surjectivity
condition is a sufficient condition for the displayed theorem, not a
necessary consequence of its hypotheses.

\subsection{ARCH-134-01: square-zero unitization of an exact scalar embedding}

\paragraph{Statement.}
Let $k$ be a field.  Suppose $\mathcal N$ is a complete and cocomplete
$k$-linear abelian category with enough projectives, and suppose
\[
 \varphi_{\mathcal N}\dashv\kappa_{\mathcal N}\dashv\psi_{\mathcal N},
 \qquad \kappa_{\mathcal N}:\operatorname{Vect}_k\longrightarrow\mathcal N
\]
is exact and fully faithful.  Let $\mathcal A$ be the category with objects
\[
       (V,X,h),\qquad V\in\operatorname{Vect}_k,
       \quad X\in\mathcal N,\quad h:\kappa_{\mathcal N}V\longrightarrow X,
\]
and morphisms $(f,g)$ satisfying $gh=h'\kappa_{\mathcal N}(f)$.
Then $\mathcal A$ is complete and cocomplete abelian with enough
projectives and carries a symmetric monoidal, colimit-preserving
$k$-bilinear tensor with unit
\[
                  \mathbf1=(k,\kappa_{\mathcal N}k,1).
\]
Its scalar functor and its two adjoints are
\[
 \kappa_{\mathcal A}U=(U,\kappa_{\mathcal N}U,1),\qquad
 \varphi_{\mathcal A}(V,X,h)=\varphi_{\mathcal N}X,\qquad
 \psi_{\mathcal A}(V,X,h)=V.
 \tag{ARCH134.1}
\]
The scalar functor is exact and strong symmetric monoidal.  Every
$\kappa_{\mathcal A}U$ is projective.  The functor
\[
       i:\mathcal N\longrightarrow\mathcal A,
                    \qquad i(X)=(0,X,0)
 \tag{ARCH134.2}
\]
is exact, preserves coproducts and projectives, and satisfies
$\varphi_{\mathcal A}i=\varphi_{\mathcal N}$.

\paragraph{Proof.}
Write $U\odot X$ for the scalar copower in $\mathcal N$,
characterized by
\[
 \operatorname{Hom}_{\mathcal N}(U\odot X,Y)
 =\operatorname{Hom}_k\bigl(U,\operatorname{Hom}_{\mathcal N}(X,Y)\bigr).
 \tag{ARCH134.3}
\]
It exists by taking a coproduct after choosing a basis of $U$.
The universal property makes it canonical and gives the coherent
isomorphisms $(U\otimes V)\odot X=U\odot(V\odot X)$ and
$k\odot X=X$.  It preserves colimits in both variables; in the second
variable this follows from its description as a coproduct of copies.
Since $\kappa_{\mathcal N}$ preserves coproducts, its $k$-linearity gives
canonical coherent identifications
\[
       U\odot\kappa_{\mathcal N}V
                     \simeq\kappa_{\mathcal N}(U\otimes_k V).
 \tag{ARCH134.4}
\]
These identifications also follow directly from (ARCH134.3) and the
right adjoint of $\kappa_{\mathcal N}$.

Limits and colimits in $\mathcal A$ are computed componentwise, using
preservation of both by $\kappa_{\mathcal N}$.  Kernels and cokernels
are therefore computed componentwise, and exactness of
$\kappa_{\mathcal N}$ supplies the structure maps.  The comparison
from coimage to image is an isomorphism in both components.  Thus
$\mathcal A$ is abelian, and its two component functors are exact.

For $(V,X,h)$ and $(W,Y,g)$ put
\[
 (V,X,h)\otimes_{\mathcal A}(W,Y,g)
  =\left(V\otimes W,
     (V\odot Y)\amalg_{\kappa_{\mathcal N}(V\otimes W)}(W\odot X),
     \bar h\right).
 \tag{ARCH134.5}
\]
The two pushout legs are $V\odot g$ and $W\odot h$, after (ARCH134.4)
and the scalar symmetry.  The structure map $\bar h$ is their common
image in the pushout.  Functoriality of the pushout gives an additive
$k$-bilinear bifunctor.

Here is an explicit coherence verification.  For a finite list
$(V_j,X_j,h_j)$, its iterated tensor has scalar part
$V=\bigotimes_jV_j$ and $\mathcal N$-part equal to the colimit of the
star diagram with central object $\kappa_{\mathcal N}V$ and one outer
object
\[
                 \left(\bigotimes_{j\ne a}V_j\right)\odot X_a
\]
for each $a$, the corresponding arrow being the copower of $h_a$.
For a one-term list this colimit is $X_1$; for a two-term list it is
(ARCH134.5).  Distributing scalar copowers over pushouts and using
(ARCH134.4) proves by induction that every parenthesization gives this
same colimit.  Define each associator through this common diagram.
All maps in the pentagon agree on the central object and every outer
object; the colimit universal property makes the pentagon commute.
Permutation of the labels and scalar factors defines the symmetry,
and the same argument verifies the hexagon.  Tensoring with
$(k,\kappa_{\mathcal N}k,1)$ gives the pushout
$X\amalg_{\kappa_{\mathcal N}V}\kappa_{\mathcal N}V=X$;
these identifications supply the unit and verify the triangle.
This establishes all symmetric monoidal axioms.  Formula (ARCH134.5)
preserves colimits in either argument, because scalar tensor and
copowers do, and colimits commute with the indicated pushout.
In particular the tensor is right exact.

A map from $(V,X,h)$ to $\kappa_{\mathcal A}U$ is uniquely determined
by its map $g:X\to\kappa_{\mathcal N}U$: full faithfulness of
$\kappa_{\mathcal N}$ supplies the unique $f$ with
$\kappa_{\mathcal N}f=gh$.  Hence these maps are naturally
$\operatorname{Hom}_k(\varphi_{\mathcal N}X,U)$.
A map in the other direction is uniquely determined by $f:U\to V$,
its second component being $h\kappa_{\mathcal N}f$.
This proves the two adjunctions in (ARCH134.1).  Exactness is
componentwise.  The pushout defining the tensor of two scalar objects
has both legs invertible, giving the canonical strong monoidal
identification $\kappa_{\mathcal A}U\otimes\kappa_{\mathcal A}W
=\kappa_{\mathcal A}(U\otimes W)$, with coherence already checked.
Thus this is indeed the scalar copower of the displayed tensor unit.

For $U\in\operatorname{Vect}_k$ and projective $P\in\mathcal N$, put
\[
            P(U,P)=(U,P\oplus\kappa_{\mathcal N}U,\mathrm{incl}).
\]
Its Hom functor is
\[
 \operatorname{Hom}_{\mathcal A}(P(U,P),(V,X,h))
       =\operatorname{Hom}_k(U,V)\times\operatorname{Hom}_{\mathcal N}(P,X).
\]
Both factors are exact, so this object is projective.  Given an object
$(V,X,h)$, choose a projective epimorphism $P\to X$; the map
$P(V,P)\to(V,X,h)$ whose second component is $(P\to X,h)$ is epic
in both components.  There are enough projectives.  Taking $P=0$
shows that every $\kappa_{\mathcal A}U$ is projective.  Taking $U=0$
shows that $i$ takes projectives to projectives.  Its other assertions
are immediate from the componentwise formulas.  This proves the
statement.\hfill$\square$

\subsection{ARCH-134-02: a full tensor-and-interval counterexample to higher surjectivity}

\paragraph{Statement.}
For every field $k$ there exists a complete and cocomplete $k$-linear
abelian symmetric monoidal category $\mathcal M$ with enough projectives
satisfying all the scalar, normalization, tensor and interval conditions
of cote~134-8, physical pp.~48--49, footers~47--48, such that the
countable comparison $b_2$ fails to be surjective.  Its tensor preserves
all colimits in each variable, and its interval and interval quotient
are flat.  The failure is explicitly
\[
                      0\longrightarrow E(k^{(\mathbb N)})\ne0,
        \qquad E(V)=V^{**}/V.
 \tag{ARCH134.6}
\]

\paragraph{Proof.}
Use the category $\mathcal N$ of ALPOGE\_BRIDGES F-SM-34.  We recall
only the precise inherited data used here.  An object is
\[
 (X_0,X_1,b:X_0\to X_1,a:E(X_1)\to X_0),\qquad
              ba=0,\quad aE(b)=0,
\]
with scalar embedding $\kappa_{\mathcal N}V=(V,0,0,0)$ and adjoints
$\varphi_{\mathcal N}X=\operatorname{coker}a$,
$\psi_{\mathcal N}X=\ker b$.  It is complete and cocomplete abelian
with enough projectives, and the scalar embedding is exact and fully
faithful.  F-SM-34, equations (SM34.6)--(SM34.8), supplies projective
resolutions giving
\[
 L_{2r}\varphi_{\mathcal N}(\kappa_{\mathcal N}V)=E^r(V),
       \qquad L_{2r+1}\varphi_{\mathcal N}(\kappa_{\mathcal N}V)=0.
 \tag{ARCH134.7}
\]
In particular $E(k)=0$, whereas $E(k^{(\mathbb N)})\ne0$.
These are inherited proved inputs, not claimed as new findings.

Apply ARCH-134-01 and let
\[
                    \mathcal M=\operatorname{Ch}_{\ge0}(\mathcal A).
\]
Give it the ordinary homological tensor formed from
$\otimes_{\mathcal A}$, with finite sums on each total diagonal and
Koszul signs.  The coherence proved above yields its symmetric
monoidal coherence.  Its unit is $\kappa_{\mathcal A}k[0]$.
Products, coproducts, kernels and cokernels are degreewise.  Standard
projective disks and degree-zero projectives show that it has enough
projectives.  More explicitly, for projective $P\in\mathcal A$,
$P[0]$ represents $\operatorname{Hom}_{\mathcal A}(P,C_0)$, and
$D_n(P)$ represents $\operatorname{Hom}_{\mathcal A}(P,C_n)$ for
$n\ge1$.  Choosing projective epimorphisms onto all $C_n$ gives an
epimorphism from a coproduct of these objects.  Each fixed degree uses
only two disk terms; exactness of arbitrary coproducts is not used.

The scalar adjunctions are
\[
 \varphi_{\mathcal M}=\varphi_{\mathcal A}H_0
       \dashv\kappa_{\mathcal M}=(\kappa_{\mathcal A}-)[0]
       \dashv\psi_{\mathcal M}=\psi_{\mathcal A}\operatorname{ev}_0.
 \tag{ARCH134.8}
\]
Indeed a chain map to a degree-zero object factors through $H_0$,
and a chain map from one is determined in degree zero.  The middle
functor is exact and strong symmetric monoidal.  Its value on every
scalar module is projective in $\mathcal M$, so the canonical
normalization is already
\[
                 L\varphi_{\mathcal M}(\kappa_{\mathcal M}U)=U[0]
                       \qquad(U\in\operatorname{Vect}_k).
 \tag{ARCH134.9}
\]
In particular condition (c) of the sheet holds with its actual counit.

We will use the natural identification
\[
       L\varphi_{\mathcal M}(C)
                    \simeq\mathbb L\varphi_{\mathcal A}(C),
 \tag{ARCH134.10}
\]
where the right-hand side regards $C$ as an $\mathcal A$-complex.
The finite-diagonal proof is included to exclude an AB4 assumption.
Resolve $C$ by projective objects $R_r$ of the chain category.  Each
$R_r$ is a retract of a coproduct of projective disks and a projective
in degree zero, so it has projective terms and projective $H_0$, with
all positive homology zero.  The first-quadrant bicomplex $(R_r)_s$
has a total complex of projectives.  Its resolution augmentation to
$C$ is a quasi-isomorphism, because its augmented rows in the
resolution direction are exact.  Its augmentation in the chain
direction to $H_0(R_\bullet)$ is also a quasi-isomorphism, because the
positive chain homology of each $R_r$ vanishes.  Both statements use
only finite sums on each diagonal.  Thus $H_0(R_\bullet)$ is a
projective complex representing $C$, and applying
$\varphi_{\mathcal A}$ computes both sides of (ARCH134.10).
Both $\operatorname{Tot}R$ and $H_0(R_\bullet)$ are bounded-below
complexes of projectives.  Their quasi-isomorphism is therefore a
chain-homotopy equivalence, which is preserved by the additive functor
$\varphi_{\mathcal A}$.  This makes the derived comparison explicit.
The augmentations show compatibility with the scalar counits.

Put $J=D_1(\mathbf1_{\mathcal A})$, and map the tensor unit of
$\mathcal M$ into its degree-zero term.  For every $C$,
\[
 (J\otimes C)_n=C_n\oplus C_{n-1},\qquad
 d(x,y)=(d_Cx+y,-d_Cy),\qquad \alpha_C(x)=(x,0).
\]
The natural map $\alpha_C$ is degreewise split monic.  The map
$H(x,y)=(0,x)$ satisfies $dH+Hd=1$.  Therefore
$J\otimes C$ is contractible and (ARCH134.10) makes its
$L\varphi_{\mathcal M}$ zero in every degree, including degree zero.
This is exactly the full condition (d).  Both $J\otimes-$ and
$(J/\mathbf1)\otimes-$ are exact: the former is the displayed finite
cone construction and the latter is shift by one.  This also proves
the optional flatness assertion.

For each $s\in\mathbb N$ take
\[
                        X_s=i(\kappa_{\mathcal N}k)[0].
\]
Since $i$ is exact and carries projectives to projectives, projective
resolutions in $\mathcal N$ show that
\[
 L_n\varphi_{\mathcal M}(iY[0])
                         =L_n\varphi_{\mathcal N}(Y).
 \tag{ARCH134.11}
\]
Put $V=k^{(\mathbb N)}$.  Preservation of coproducts by $i$ and
$\kappa_{\mathcal N}$ identifies
$\bigoplus_sX_s=i(\kappa_{\mathcal N}V)[0]$.  Formula (ARCH134.7)
now gives zero for every summand on the left side of $b_2$ and
$E(V)\ne0$ for its right side.  This proves (ARCH134.6) and all
assertions.\hfill$\square$

\subsection{ARCH-134-03: its localization is still the connective derived category}

\paragraph{Statement.}
In the example ARCH-134-02, set
$F=L\varphi_{\mathcal M}$ and $W=F^{-1}(\mathrm{iso})$.
Then the actual induced functor is an equivalence
\[
                         \mathcal M[W^{-1}]\simeq D_{\ge0}(k).
 \tag{ARCH134.12}
\]
Nevertheless $W$ is not closed under countable coproducts.  Thus neither
higher coproduct surjectivity nor closure of $W$ under coproducts is
necessary for the ordinary localization conclusion in the sheet.

\paragraph{Proof.}
The exact scalar functor $\kappa_{\mathcal A}$ is right adjoint to
$\varphi_{\mathcal A}$ and sends every scalar module to a projective
object.  On connective derived categories we therefore have the
adjunction
\[
 \mathbb L\varphi_{\mathcal A}:D_{\ge0}(\mathcal A)
       \rightleftarrows D_{\ge0}(k):\kappa_{\mathcal A}.
 \tag{ARCH134.13}
\]
For clarity, the adjunction follows by representing the source by a
bounded-below projective complex: applying $\varphi_{\mathcal A}$
gives a complex of projective scalar modules, and the ordinary
termwise adjunction computes the two derived Hom groups.  The usual
homotopy lifting argument for bounded-below complexes of projectives
requires no AB4 axiom.

A connective scalar complex $C$ is sent by $\kappa_{\mathcal A}$ to
a complex of projectives.  Consequently its canonical counit is
represented literally by
\[
       \varphi_{\mathcal A}\kappa_{\mathcal A}(C)=C.
\]
Hence the right adjoint in (ARCH134.13) is fully faithful.  An
adjunction with invertible counit is a localization at the maps
inverted by its left adjoint: its units belong to that class by the
triangle identity, and after formally inverting those units the two
functors become mutually inverse.  This elementary argument proves
that localization assertion directly.

By (ARCH134.10), the functor $F$ on the chain category factors through
its localization at quasi-isomorphisms as the left functor of
(ARCH134.13).  Localizing first at quasi-isomorphisms, then at the
remaining maps inverted by that functor, proves (ARCH134.12).

Here is the explicit failure of closure under coproducts.  In
$\mathcal A$ there is the morphism
\[
 u_U:i(\kappa_{\mathcal N}U)\longrightarrow\kappa_{\mathcal A}U
\]
whose components are $0\to U$ and $1_{\kappa_{\mathcal N}U}$.
For $U=k$, (ARCH134.7) and (ARCH134.9) show that $u_k[0]$ is in
$W$, with the identity on degree-zero homology.  Its countable
coproduct is $u_V[0]$, $V=k^{(\mathbb N)}$.  The source of its
$F$-image has nonzero degree-two homology $E(V)$, whereas its target
has zero degree-two homology.  Therefore $u_V[0]\notin W$.
This proves the second assertion.\hfill$\square$

\subsection{What this does and does not decide}
ARCH-134-02 refutes the universal higher-surjectivity assertion asked
for in task~2, under the full original predicates.  It strictly
advances F-SM-43: in that earlier non-AB4 example the derived
coproduct comparison remained an isomorphism, whereas here it does
not.  ARCH-134-03 checks the localization conclusion in the new
counterexample and finds it true.  Thus this example is not a
counterexample to Grothendieck's \emph{Thm(?)} itself.  A proof of that
theorem in complete abstract generality must avoid imposing higher
coproduct compatibility as a consequence of the handwritten axioms.

The example does not claim a small set of finitely presented
projective generators.  In particular it does not meet the additional
ringoid corollary by an asserted presentation, and it does not decide
SM25.2 for the raw ordinary right adjoint $G$.  That separate gate
remains the simultaneous derived acyclicity of the products of
scalar-coinduced projectives written there.  The construction here
supplies no violation of that Tor assertion and no proof of it.

\paragraph{Geometric language and limits.}
The scalar adjoints are the linear analogue of taking homology and
constant coefficients.  The result separates the existence of a
connective homology-type realization from preservation of infinite
sums before localization.  It is a statement about Grothendieck's
linearization framework on the cited sheet; it does not assert a new
property of schemes over $\mathbf Z$, $\lambda$-rings, algebraic
$K$-theory, motives, or $\ell$-adic weights.  No EGA, SGA, or FGA
statement is used as an input to this example, so no fictitious
locator in those works is supplied.
